\centerline{\bf \Huge Налоговые вычеты I}

\begin{flushright}
        \scriptsize{}
\end{flushright}

\problem{1}
Докажите, что для любого простого $p$ существует решение сравнения

$$
\left(x^2-3\right)\left(x^2-5\right)\left(x^2-15\right) \equiv 0 \quad(\bmod p) 
$$

\problem{2}
Решите в целых числах уравнение $x^2=y^3+7$.

\textit{Указание: прибавьте к обеим частям 1 и правую часть разложите на множители, после чего рассмотрите обе части по модулю $p$ делителя числа $x^2+1$}

\problem{3}
Числа $a^2+5 a+1$ и $2 a^2-a+2$ не делятся на простое число $p$ ни при каких целых a. Докажите, что для некоторого натурального $n$ число $n^2+3 n+11$ делится на $p$.

\textit{Указание: воспользуйтесь знанием о том, что $ax^2 + bx + c \equiv 0 \ (mod \ p)$ имеет решение тогда и только тогда, когда $\left(\dfrac{D}{p}\right) = \left(\dfrac{b^2-4ac}{p}\right) = 1$}

\problem{4}
Пусть $p=4 k-1$ --- простое число. Докажите, что если сравнение $x^2 \equiv a \ (mod \ p)$ имеет решение, то $x \equiv \pm a^k \ (mod \ p)$.

\textit{Указание: воспользуйтесь критерием Эйлера $\left(\dfrac{a}{p}\right) \equiv a^{\dfrac{p-1}{2}}$}